\begin{figure}[htbp]
    \centering
    \begin{tikzpicture}[scale=1.5]
      % 定义三个圆的中心
      \coordinate (A) at (0,0);
      \coordinate (B) at (1.5,0);
      \coordinate (C) at (0.75,1.3);
  
      % 画三个圆
      \draw[circle] (A) circle (1.9);
      \draw[circle] (B) circle (1.9);
      \draw[circle] (C) circle (1.9);
  
      % 标注圆的名称
      \node (finite) at ($(A)+(-0.6,-0.5)$) [above left] {\heiti{有限}};
      \node (finite_ref) at (finite.south) [below] {\textbf{finite}};
      \node (discrete) at ($(B)+(0.6,-0.5)$) [above right] {\heiti{离散}};
      \node (discrete_ref) at (discrete.south) [below] {\textbf{discrete}};
      \node (algorithm_relatedness) at ([yshift=12pt]$(C)+(0,0.7)$) [above] {\heiti{算法相关}};
      \node at (algorithm_relatedness.south) [below] {\textbf{algorithm-relevant}};
  
      % 圆之间的交集文字
      \node (PAC) at ($(A)!0.5!(B)+(0,-0.6)$) [below] {PAC\ VCD\ SC};
      \node (NTK) at ($(B)!0.5!(C)+(0.4,0.7)$) [right] {NTK};
      \node (EA) at ([xshift=-1pt,yshift=-10pt]NTK) [right] {EA};
      \node at ($(C)!0.5!(A)+(-0.2,0.6)$) [left] {MF\ GF};
  
      % 三圆交集文字
      \node (FLT) at ($(C)+(0,-0.6)$) {FLT};
      \node (FLT_ref) at (FLT.south) [below] {\cite{DBLP:conf/focs/Allen-ZhuL21}};
  
      % Caption
      \coordinate (caption1) at ($(A)+(-6.5,2)$);
      \node at (caption1) [right] {均值场理论（MF）: \cite{DBLP:conf/focs/Allen-ZhuL21}};
      \coordinate (caption2) at ([yshift=-12pt]caption1);
      \node at (caption2) [right] {梯度流（GF）: \cite{DBLP:conf/focs/Allen-ZhuL21}};
      \coordinate (caption3) at ([yshift=-12pt]caption2);
      \node at (caption3) [right] {神经正切核（NTK）: \cite{DBLP:conf/focs/Allen-ZhuL21}};
      \coordinate (caption4) at ([yshift=-12pt]caption3);
      \node at (caption4) [right] {概率近似正确（PAC）: \cite{DBLP:conf/focs/Allen-ZhuL21}};
      \coordinate (caption5) at ([yshift=-12pt]caption4);
      \node at (caption5) [right] {可表示性理论（EA）: \cite{DBLP:conf/focs/Allen-ZhuL21}};
      \coordinate (caption6) at ([yshift=-12pt]caption5);
      \node at (caption6) [right] {VC 维理论（VCD）: \cite{DBLP:conf/focs/Allen-ZhuL21}};
      \coordinate (caption7) at ([yshift=-12pt]caption6);
      \node at (caption7) [right] {样本复杂度（SC）: \cite{DBLP:conf/focs/Allen-ZhuL21}};
  \end{tikzpicture}
  \caption{学习理论分类图。}
  \label{fig-learning-theory-taxonomy}
  \end{figure}